\section*{ЗАКЛЮЧЕНИЕ}
\addcontentsline{toc}{section}{ЗАКЛЮЧЕНИЕ}

Разработанная программная система оптимизирующего компилятора Common Lisp представляет собой полностью функциональный транслятор исходного кода на языке Common Lisp в байт-код, реализующий все основные возможности языка, а также виртуальную машину для выполнения исходных программ.

Основные результаты работы:

\begin{enumerate}
\item Изучен процесс компиляции языков программирования и функциональных языков в частности. Сформулированы основные стадии компиляции.
\item Изучены методы оптимизации компиляторов. Подобран список оптимизаций, которые возможно реализовать при компиляции Common Lisp.
\item Спроектирован компилятор с оптимизациями. Разработана архитектура компилятора и программный интерфейс. Составлен список реализуемых оптимизаций.
\item Реализована компиляция S-выражения в промежуточную форму, состоящую из простых операций и сохраняющую смысл исходного выражения, а также сбор информации о функциях для их дальнейшей оптимизации.
\item Реализована генерация кода ассемблера. Для этого происходит обход дерева промежуточной формы в глубины и для каждой операции определённым образом генерируется набор линейных инструкций ассемблера, состоящих из мнемоник и операндов, а также меток для ветвления и функций.
\item Реализована генерация байт-кода. Происходит в 2 прохода. В первый проход в каждой инструкции мнемоника заменяется на соответствующий опкод, константы собираются в список без дубликатов, а обращения к этим константам содержат индексы в этом списке, для меток запоминается их адрес в байт-коде, инструкции переходов и вызова функции сохраняются для дальнейшей замены меток. Во втором проходе все упоминания меток заменяются на относительные адреса, полученные в первом проходе.
\item Реализована оптимизация промежуточной формы. Промежуточное дерево, получившееся в результате статического анализа, преобразуется посредством применения оптимизаций. Происходит полный обход дерева в глубину, где для разных операций предусмотрены различные виды оптимизаций, и вызываются соответствующие функции.
\end{enumerate}

Все требования, объявленные в техническом задании, были полностью реализованы, все задачи, поставленные в начале разработки проекта, были также решены.

Готовый рабочий проект представлен в виде программной системы компилятора Common Lisp. Исходный код находится в публичном доступе и внедрён в os-pi.
